% ============================================================================
%  Appendix: neural teleportation is an exact isomorphism; the measured logit
%  drift is floating-point roundoff.
%
%  Added 2026-09-07. This appendix REPLACES the earlier "teleportation is
%  function-approximate" diagnosis, which attributed the measured drift to
%  un-migrated BatchNorm running statistics. That diagnosis was a conjecture
%  (docs/Final-twist/km-notes.md, 2026-05-01, "Most likely cause: ...") and it
%  is wrong: the library does not migrate the running statistics because it does
%  not need to. See scripts/verify_teleportation_exactness.py.
% ============================================================================

\section{Neural teleportation is exact: the measured drift is roundoff}
\label{app:teleport-exact}

That neural teleportation preserves the function exactly is a theorem, not a
finding of this paper: an isomorphism of neural networks $\tau:(W,f)\to(V,g)$
satisfies $\Psi(W,f)=\Psi(V,g)$ \citep[Thm.~4.13]{armenta2021representation}, and
a per-neuron change of basis is such an isomorphism. This appendix does two
things that the theorem does not. Section~\ref{app:teleport-why} records why
batch normalisation is inside the theorem's scope rather than an exception to it,
since that is the point at which the construction is most often doubted.
Sections~\ref{app:teleport-eps}--\ref{app:teleport-tf32} then verify that our
\emph{software} realises the identity, and calibrate the arithmetic floor below
which no drift figure in this paper can be resolved --- a question about the
implementation, on which a theorem is silent.

\subsection{Why the transform is exact, batch normalisation included}
\label{app:teleport-why}

A teleportation assigns a nonzero scalar $\tau_q$ to each hidden neuron $q$ and
conjugates the parameters, $W^{(\ell)}\mapsto T_\ell W^{(\ell)}T_{\ell-1}^{-1}$
and $b^{(\ell)}\mapsto T_\ell b^{(\ell)}$ with $T_\ell=\mathrm{diag}(\tau)$. The
group acts on the activations as well, by
\begin{equation}
(\tau\cdot f)_v(z)\;=\;\tau_v\,f_v\!\left(\frac{z}{\tau_v}\right)
\label{eq:cob-activation}
\end{equation}
at each hidden vertex $v$ \citep[Eq.~2]{armenta2021representation}. For a
positively homogeneous activation --- ReLU with $\tau_v>0$ --- this leaves $f_v$
unchanged, the $T_\ell$ telescope through the layer stack, and the function is
preserved.

Batch normalisation is the case that looks like an obstruction, and it is worth
being explicit about why it is not. In evaluation
mode a BN channel is the affine
map $z\mapsto\gamma(z-\mu)/\sqrt{\sigma^2+\epsilon}+\beta$, which is
\emph{not} positively homogeneous: the running mean $\mu$ and the shift $\beta$
break $\varphi(\tau z)=\tau\varphi(z)$. One repair is to migrate the running
statistics, $\mu\mapsto\tau\mu$ and $\sigma^2\mapsto\tau^2\sigma^2$; note that
this repair is only exact if the numerical guard $\epsilon$ is rescaled too,
since $\sqrt{\tau^2\sigma^2+\epsilon}\neq\tau\sqrt{\sigma^2+\epsilon}$.

None of that is necessary, because the framework already prescribes the answer.
Two facts settle it. First, in evaluation mode $\mu$ and $\sigma^2$ are not
batch-dependent quantities but ordinary weights of the network
\citep[Remark~5.4]{armenta2021representation}, so a normalisation layer is a pair
of affine vertices like any other and Theorem~4.13 applies to it unchanged.
Second, the action on its activation is given by
Equation~\eqref{eq:cob-activation}, and that is precisely what the
\texttt{neuralteleportation} implementation computes: it divides the incoming
activation by the incoming change of basis, applies the unmodified
normalisation, and carries the outgoing basis on the affine parameters,
\[
z \;\xrightarrow{\;\text{incoming COB}\;}\; \tau_{\mathrm{in}} z
\;\xrightarrow{\;\div\,\tau_{\mathrm{in}}\;}\; z
\;\xrightarrow{\;\mathrm{BN}_{\mu,\sigma^2,\epsilon}\;}\; \mathrm{BN}(z)
\;\xrightarrow{\;\gamma,\,\beta\;\times\;\tau_{\mathrm{out}}\;}\;
\tau_{\mathrm{out}}\,\mathrm{BN}(z).
\]
which is Equation~\eqref{eq:cob-activation} with $f_v=\mathrm{BN}$. The division
restores the original pre-normalisation activation exactly, so the original
$\mu$, $\sigma^2$ and $\epsilon$ remain correct, and the composite is
$\tau_{\mathrm{out}}\mathrm{BN}(z)$ --- with no condition on $\epsilon$, and no
migration of the running statistics. Teleportation of a batch-normalised network
in evaluation mode is therefore function-exact, and satisfies the hypotheses of
Theorem~\ref{thm:maxinv} and Lemma~\ref{lem:quiver-inv} exactly.

We note one boundary of the framework that bears on this paper elsewhere:
\citet{armenta2021representation} observe that average and global-average pooling
sit inside it without qualification, whereas max-pooling ``break[s] the algebraic
structure'' --- the same tie-dependence that Theorem~\ref{thm:weff} excludes by
hypothesis.

\subsection{The measured drift scales with machine epsilon}
\label{app:teleport-eps}

Exactness in real arithmetic does not mean a zero measurement in floating point.
The round trip above --- multiply by $\tau_{\mathrm{in}}$ inside the preceding
convolution, then divide by $\tau_{\mathrm{in}}$ inside the normalisation --- is
exact over $\R$ but not over floats: the convolution is recomputed with rescaled
weights, so its partial sums have different magnitudes and round differently.
The drift is therefore $O(\varepsilon_{\mathrm{mach}})$ and, decisively,
\emph{linear in the working precision}.

That is a falsifiable prediction, and it separates the two hypotheses cleanly. If
the transform genuinely changed the function --- if the running statistics really
were stale --- the discrepancy would be a property of the weights and would be
essentially \emph{independent} of working precision, giving a
\texttt{fp32}/\texttt{fp64} ratio near $1$. If it is roundoff, the ratio must be
near $\varepsilon_{32}/\varepsilon_{64}=2^{29}\approx5.4\times10^{8}$.

We ran the identical change-of-basis draw (same seed, same COB vector) at both
precisions on all three architectures, with the library's three internal
\texttt{float32} down-casts patched out so that the \texttt{fp64} arm is
genuinely double precision.\footnote{The library hard-casts to single precision
in \texttt{COBForwardMixin.forward}, in \texttt{layers/merge.py:Add.forward} (the
residual skip connection), and via a \texttt{float32} tracing input in
\texttt{NeuralTeleportationModel.\_\_init\_\_}. Without these three patches the
\texttt{fp64} arm silently retains a single-precision floor and the ratio
saturates near $10$.} Table~\ref{tab:teleport-exactness} reports the result.

\begin{table}[t]\centering\small
\caption{Teleportation is exact; the drift is roundoff. Maximum absolute logit
discrepancy $\max|\Net-\Psi_\tau(W\!,f)|$ between a network and its teleported image, for the
\emph{same} change-of-basis draw evaluated at two working precisions ($5$ seeds,
$8$ inputs, \texttt{cob\_range}$=1$, CPU, evaluation mode). The \texttt{fp64}
discrepancy sits at the double-precision floor ($\sim10^{-15}$ relative), and the
ratio between the two columns matches
$\varepsilon_{32}/\varepsilon_{64}=5.37\times10^{8}$ --- the signature of pure
roundoff. A genuine change of function would give a ratio near $1$.}
\label{tab:teleport-exactness}
\smallskip
\begin{tabular}{@{}lccc@{}}\toprule
Architecture & \texttt{fp32} drift & \texttt{fp64} drift (relative) & \texttt{fp32}/\texttt{fp64} ratio \\ \midrule
ResNet-152   & $4.2$--$6.1\times10^{-6}$ & $8.2\times10^{-15}$--$1.0\times10^{-14}$ \; ($1.4$--$1.7\times10^{-15}$) & $4.2$--$6.7\times10^{8}$ \\
DenseNet-121 & $6.2$--$10.5\times10^{-6}$ & $1.2$--$2.3\times10^{-14}$ \; ($1.6$--$3.2\times10^{-15}$) & $3.3$--$7.9\times10^{8}$ \\
GoogLeNet    & $5.7$--$8.6\times10^{-6}$ & $1.4$--$1.7\times10^{-14}$ \; ($1.9$--$2.2\times10^{-15}$) & $3.4$--$5.7\times10^{8}$ \\
\bottomrule\end{tabular}\end{table}

The ratio is $3.3$--$7.9\times10^{8}$ across all fifteen runs, straddling the
predicted $5.37\times10^{8}$; the \texttt{fp64} residual is $1.4$--$3.2\times
10^{-15}$ relative, i.e.\ at the double-precision floor. The measured drift is
roundoff, and nothing else.

\subsection{The knowledge matrix itself, not merely the logits}
\label{app:teleport-km}

The logit test above certifies that the \emph{function} is preserved. The
knowledge matrix is a strictly finer object --- it records the germ, which
Theorem~\ref{thm:nogo} shows the logits do not control --- so it deserves its own
check. Writing $\Psi_c$ for class $c$'s component of $\Net$, we recomputed
knowledge-matrix rows directly, as
$\M_{c,:}(x)=[\,\nabla_x\Psi_c(x)\odot x \mid \Psi_c(x)-\nabla_x\Psi_c(x)^{\!\top}x\,]$
by one vector--Jacobian product per class, for the original and the teleported
network, at both precisions ($12$ class rows, \texttt{cob\_range}$=1$).

At double precision the knowledge-matrix drift is
$0.8$--$2.5\times10^{-15}$ relative across all six runs --- the same
machine-epsilon floor as the logits. \emph{The knowledge matrix is invariant under teleportation to double
precision}, on ResNet-152, DenseNet-121 and GoogLeNet alike, exactly as
Lemma~\ref{lem:quiver-inv} requires.

Two scope notes, so this test is not read for more than it establishes. First,
the bias column here is \emph{derived} as $\Psi_c(x)-\nabla_x\Psi_c(x)^{\!\top}x$
rather than accumulated independently, so the row-sum identity holds by
construction and is not itself under test; what the comparison tests --- and what
Lemma~\ref{lem:quiver-inv} is about --- is that the Jacobian block
$\nabla_x\Psi(x)\,\mathrm{diag}(x)$ is identical between the original and
teleported networks. Second, this route computes the knowledge matrix by
vector--Jacobian products on the teleported model directly, and therefore
deliberately bypasses the library's serialised change-of-basis loading path. That
path has a separate, genuine defect --- a knowledge matrix loaded through it
violates $\M\mathbf 1=\Net$ structurally --- which is why the measured-drift
column of the earlier teleportation run was withdrawn, and which the present
result does not repair. The two findings are about different code, and both
stand: the transform is exact, and the loader is buggy.

One detail is worth recording rather than smoothing over. The
\texttt{fp32}/\texttt{fp64} ratio for the knowledge matrix is
$7\times10^{10}$--$2.4\times10^{12}$,
\emph{larger} than the $5.4\times10^{8}$ epsilon ratio that the logits obey. The
excess is not evidence of a function change --- a function change would push the
ratio \emph{towards} $1$, not away from it. It is the ill-conditioning of the
change of basis (Section~\ref{app:teleport-tf32}) acting on a differentiated
quantity: the backward pass propagates through the same $\tau$ factors as the
forward pass, so single precision loses accuracy super-linearly where double
precision still has the headroom to stay at its floor. The decisive fact is the
\texttt{fp64} column, which is at machine epsilon.

\subsection{Why the cluster run measured $10^{-2}$ rather than $10^{-6}$}
\label{app:teleport-tf32}

The single-precision drift above is $\sim\!10^{-6}$, three orders of magnitude
\emph{below} the $10^{-3}$ acceptance threshold that the original
equivalence check applied --- yet that check reported every teleportation
exceeding it, at $3.1\times10^{-2}$ to $1.0\times10^{-1}$ on ResNet-152. The gap
is the accelerator's arithmetic, not the mathematics.

The teleportation run ran on H100 GPUs (\texttt{job\_teleportation.sh} requests
\texttt{-{}-gpus=h100:1}). PyTorch enables TF32 for cuDNN convolutions by
default, and this repository never disables it. TF32 carries a $10$-bit
mantissa, so $\varepsilon_{\mathrm{TF32}}=2^{-11}=4.9\times10^{-4}$, coarser
than single precision by a factor of $8192$. Since the drift is linear in
$\varepsilon$, scaling the measured CPU \texttt{fp32} drift by that factor
predicts a TF32 drift of $3.5$--$5.0\times10^{-2}$ on ResNet-152, against
$3.1\times10^{-2}$--$1.0\times10^{-1}$ observed. For DenseNet-121 and GoogLeNet
the prediction ($5$--$9\times10^{-2}$) overshoots the observed
$1.6$--$3.9\times10^{-2}$ by a factor of two to four, as expected --- not every
kernel in those graphs runs in TF32, and our inputs are synthetic rather than
natural images. The mechanism is established by the precision-scaling law of
Section~\ref{app:teleport-eps}; the TF32 arithmetic explains the particular
magnitude the cluster recorded.

A second, independent amplifier is the conditioning of the change of basis
itself. With \texttt{cob\_range}$=1$ the library samples $\tau$ uniformly on
$[0,2]$, an interval whose closure contains zero: over the $\sim\!2\times10^{4}$
neurons of these networks the smallest drawn $\tau$ is $\sim\!3\times10^{-5}$,
so the multiply-then-divide round trip passes through a factor of $\sim\!10^{4}$.
This is why the roundoff constant is large in absolute logit units, and it is the
same conditioning problem that forced \texttt{cob\_range}$=0.1$ for ResNet-152 in
the knowledge-matrix drift run (where the per-path product of $152$ such factors
overflowed single precision outright). Sampling $\tau$ on an interval bounded
away from zero is the appropriate fix and costs nothing: the transform is
function-preserving at any range.

\subsection{What this establishes, and what it does not}
\label{app:teleport-consequences}

\begin{enumerate}
\item \textbf{The transform satisfies the theorems' hypothesis.} Teleportation is
an exact quiver isomorphism (Lemma~\ref{lem:quiver-inv}), so the knowledge-matrix
drift under it is zero in exact arithmetic and the teleportation check covers a multi-architecture
verification of invariance on all three of the residual, dense and inception
topologies.

\item \textbf{It does not supply evidence for the cross-architecture claim.}
Invariance under a quiver isomorphism is automatic by construction, so this
result --- like the permutation check of the permutation check --- confirms that the
implementation realises an identity the theory already guarantees. The
non-automatic content of Theorem~\ref{thm:maxinv} is invariance across
\emph{different architectures} realising one function, which no
same-architecture transform can exercise (Section~\ref{sec:limitations}, L1).

\item \textbf{The three-tier structure of Proposition~\ref{prop:visdrift} applies
at tier~(i).} The visible/invisible split, the no-go theorem
(Theorem~\ref{thm:nogo}) and the conditional bound remain correct and remain
worth stating --- they are what one needs for a transform that is genuinely
approximate. Teleportation is not such a transform.

\item \textbf{The measured drift is a resolution figure, not a property of the
symmetry.} It belongs with the permutation noise floor of the permutation check --- both are
finite-precision accumulation --- and it falls to the double-precision floor by
computing in \texttt{fp64}, or on the accelerator by disabling TF32
(\texttt{torch.backends.cudnn.allow\_tf32 = False}) and sampling the change of
basis away from zero.
\end{enumerate}

\subsection{The two implementation checks in full}
\label{app:teleport-checks}

The setup, procedure and measured values of the permutation and teleportation
checks are collected here. Neither establishes a mathematical fact --- both
transforms are quiver isomorphisms, so the invariance they exercise is
guaranteed by Theorem~4.13 of \citet{armenta2021representation} --- and both are
reported for what they are: verification that the software realises the identity,
and calibration of the floor below which no drift figure in this paper can be
resolved.

% ============================================================================
% Rewritten study text (drop-in snippets). Tables referenced here live in
% ../study2_tables/ and are regenerated by scripts/make_study2_tables.py.
% ============================================================================

% ---------------------------------------------------------------------------
% STUDY 1a — replaces the "130x / identical to numerical precision" framing
% ---------------------------------------------------------------------------
\subsection{Study 1A: permutation invariance (a correctness check) (A1)}
\label{sec:study1a}
\label{sec:A1}

\paragraph{Why it matters.}
A neuron permutation is a quiver isomorphism, so knowledge-matrix
invariance under it is automatic by construction
(Proposition~\ref{prop:perm} is a corollary of quiver-isomorphism
invariance, not a hard-won discovery). This experiment is therefore a
numerical \emph{correctness check}: it confirms that the implementation
realises the identity the theory guarantees, and quantifies the float32
accumulation that separates the computed object from its exact value. It
is not the paper's primary evidence for invariance; the load-bearing,
non-automatic invariance is the cross-architecture case
(Theorem~\ref{thm:maxinv}, Study~3), whose only transform-based evidence
is the approximate teleportation of Study~1B.

\paragraph{Setup.}
Architecture: ResNet-152 (the only trio member that admits an in-place
channel permutation; DenseNet-121 and GoogLeNet have concatenation
topologies with no in-place channel swap and are covered by
teleportation instead). Transform: random permutations of the
\emph{wide-face} permutation subgroup --- the 2048-channel post-pool
faces; bottleneck interiors are not permutable in place. Inputs for the
adversarial signal: clean/adversarial pairs from the six-attack suite
(Section~\ref{sec:setup}). We measure both drifts in raw fp64 Frobenius
(KM) and $\ell_2$ (penultimate), reported \emph{relative to the same
model's adversarial-pair signal} so the comparison is unit-free.

\paragraph{Procedure.}
\begin{enumerate}
\item Compute $\M(x)$ and $h(x)$ for each input.
\item Draw a random wide-face neuron permutation $P$ and build the
isomorphic network $\tilde f = P \cdot f$ (identical function, permuted
weights).
\item Recompute $\M_{\tilde f}(x)$ and $h_{\tilde f}(x)$.
\item Record the permutation \emph{noise} $\|\M_{\tilde f}(x)-\M(x)\|_F$
and $\|h_{\tilde f}(x)-h(x)\|_2$.
\item Record the adversarial \emph{signal} $\|\M(x')-\M(x)\|_F$ and
$\|h(x')-h(x)\|_2$ over the attack pairs, and report each noise relative
to its own signal.
\end{enumerate}

\paragraph{What we measure and the claim it licenses.}
The signal-to-noise ratio per representation. The claim it licenses is
narrow and correct: the knowledge matrix is permutation-invariant up to
float32 accumulation, whereas a statistic of the permuted penultimate
features moves by an amount of the same order as the adversarial signal
it is meant to register --- exactly as Theorem~\ref{thm:maxinv} requires.
It does \emph{not} license a superiority claim: a permutation-equivariant
penultimate statistic would also be invariant, and the gradient$\times$input
family shares the knowledge matrix's invariance (Theorem~\ref{thm:weff}).

\paragraph{Result.}
Measured signal-relative (Table~\ref{tab:signal-relative}), the
knowledge-matrix permutation noise is $0.1$--$0.2\%$ of its
adversarial signal in the mean (the signal sits
$7{\times}10^{2}$--$1.7{\times}10^{7}$ above the noise floor across
architectures and attacks), whereas the penultimate drift is of the
same order as its own adversarial signal (at most $3.05\times$ below it,
and \emph{above} it in 16 of 24 cells). Honesty rows we commit to: the
worst-case tail ratio on ResNet-152 (minimum signal over maximum noise)
reaches $1.14\times$ for DeepFool --- the extreme tails nearly touch on
the BN-heaviest architecture; the residual KM drift (mean $0.116$, max
$5.78$ Frobenius) is float32 accumulation, not a violation of the
theorem --- float64 spot checks reduce it to numerical zero. DenseNet-121
and GoogLeNet admit no in-place channel permutation (concatenation
topologies) and are covered by teleportation instead.

\paragraph{Reading.}
The result shows that the implementation realises the construction's
guaranteed permutation invariance to float32 accumulation, and that a
penultimate statistic does not. It does \emph{not} show that this
invariance is unique to knowledge matrices, nor that the matrix is a
better representation --- only that the identity holds where the theory
says it must, on the one trio member where the permutation is realisable
in place.

% ---------------------------------------------------------------------------
% STUDY 1b — replaces the theorem-asserted zero
% ---------------------------------------------------------------------------
\subsection{Study 1B: teleportation, under the three-tier claim structure (A2)}
\label{sec:study1b}
\label{sec:A2}

\paragraph{Why it matters.}
Teleportation is the only \emph{multi-architecture} invariance arm and the
only one that exercises a transform broader than a quiver isomorphism. But
on these BatchNorm networks it is function-\emph{approximate}, so the
invariance theorem's hypothesis does not hold and no table entry may be
asserted from it. This study is the honest reporting of that situation:
it measures what can be measured and bounds only what the violated
hypothesis still permits.

\paragraph{Setup.}
Architectures: all three trio members (ResNet-152, DenseNet-121,
GoogLeNet). Transform: neural teleportation \citep{armenta2024neural} ---
a function-preserving change-of-basis (per-neuron rescaling) realised by
the \texttt{neuralteleportation} library's COB models. Canonical run:
$T = 50$ random teleportations per architecture (seed indices
$0,\ldots,49$), evaluated on $N = 25{,}000$ ImageNet-validation samples
(the first $25{,}000$ by sorted filename), matching the compiled panel
(Table~\ref{tab:s1_table}, Section~\ref{sec:study1c}). We verify
function-equivalence per teleportation via the $10^{-3}$ logit-drift
gate.

\paragraph{Procedure.}
\begin{enumerate}
\item Compute $\M(x)$, $h(x)$, $f(x)$ for each sample.
\item Draw a random COB teleportation $\tau$ and build $\tilde f = \tau
\cdot f$, an approximate-isomorphism image.
\item Measure the logit gate $\varepsilon = \|f_\tau(x) - f(x)\|$ per
sample and test it against $10^{-3}$.
\item Compute the \emph{visible} knowledge-matrix drift
$\|(\M_\tau(x)-\M(x))\mathbf 1\|$ and verify per pair that it equals the
gate $\varepsilon$ (Proposition~\ref{prop:visdrift}(ii)).
\item Recompute the penultimate features $h_\tau(x)$ and pass the pairs
to the similarity panel of Study~1C (Section~\ref{sec:study1c}).
\end{enumerate}

\paragraph{What we measure and the claim it licenses.}
The logit gate, the visible knowledge-matrix drift (an \emph{identity}
equal to the gate by the row-sum identity, not a measurement of
invariance), and the penultimate drift that the Study~1C panel
adjudicates. The claim this licenses is exactly the three-tier structure
of Proposition~\ref{prop:visdrift}: under an exact isomorphism the drift
would be zero; under this approximate teleportation the visible component
equals the gate; the invisible component carries no a priori bound from
the gate (Theorem~\ref{thm:nogo}) and is \emph{not} estimated here.

\paragraph{Result.}
In our runs, $100/100$ teleportations exceed the $10^{-3}$ logit-drift
gate (up to $10^{-1}$ on ResNet-152, under $4\times10^{-2}$ otherwise),
so the invariance theorem's hypothesis does not hold and no table entry
may be asserted from it. We report the knowledge-matrix drift through the
visible/invisible split: the visible channel obeys the equality
$\|(\M_\tau(x)-\M(x))\mathbf 1\|=\|f_\tau(x)-f(x)\|$, which we verify per
pair against the logit gate; the invisible component is left unestimated.
The penultimate features drift substantially under teleportation (the
panel of Study~1C reports Procrustes shape distance $632$--$1471$ and
soft-matching distance $22$--$61$ across the trio,
Table~\ref{tab:s1_table}); against these the panel adjudicates which
measures quotient the drift out.

\paragraph{Reading.}
This study does \emph{not} verify invariance across the three
architectures: its transform fails the isomorphism hypothesis $100/100$.
What it shows is the visible knowledge-matrix drift pinned to the
function-approximation gate by an identity, with the invisible component
honestly flagged as unbounded and unmeasured. We credit the three-tier
framing as the correct reporting of an approximate transform, not as
evidence the matrix is invariant under one.


% Signal-relative invariance: permutation noise vs adversarial signal (raw fp64 Frobenius / L2 both sides).
\begin{table}[t]
\centering\small
\begin{tabular}{lcccc}
\toprule
\multicolumn{5}{l}{\emph{(a) Permutation-noise floors (KM Frobenius; penultimate L2)}}\\
architecture & KM mean & KM median & KM max & penult.\ mean \\
\midrule
ResNet-152 & 0.116 & 0.005137 & 5.78 & 15.05 \\
ResNet-18 & 0.0704 & 0.0107 & 0.865 & 29.32 \\
AlexNet & 1.54e-05 & 1.455e-05 & 3e-05 & 105.7 \\
VGG & 2.13e-05 & 2.053e-05 & 3.84e-05 & 57.7 \\
\midrule
\multicolumn{5}{l}{\emph{(b) Adversarial signal-to-noise (mean signal / mean noise): KM $\mid$ penultimate}}\\
attack & ResNet-152 & ResNet-18 & AlexNet & VGG \\
\midrule
FGSM & $726 \mid 0.51$ & $3257 \mid 0.63$ & $8.7\mathrm{e}6 \mid 0.75$ & $8.3\mathrm{e}6 \mid 0.69$ \\
PGD & $1029 \mid 1.36$ & $4739 \mid 1.84$ & $1.3\mathrm{e}7 \mid 1.82$ & $1.7\mathrm{e}7 \mid 3.05$ \\
CW & $845 \mid 0.48$ & $2369 \mid 0.36$ & $4.1\mathrm{e}6 \mid 0.24$ & $3.9\mathrm{e}6 \mid 0.28$ \\
DeepFool & $984 \mid 0.34$ & $1364 \mid 0.13$ & $2.9\mathrm{e}6 \mid 0.15$ & $2.3\mathrm{e}6 \mid 0.14$ \\
APGD & $803 \mid 1.02$ & $3430 \mid 1.03$ & $1.3\mathrm{e}7 \mid 1.71$ & $1.1\mathrm{e}7 \mid 1.80$ \\
Square & $914 \mid 0.48$ & $2748 \mid 0.37$ & $6.0\mathrm{e}6 \mid 0.42$ & $4.9\mathrm{e}6 \mid 0.35$ \\
\bottomrule
\end{tabular}
\caption{Permutation invariance, signal-relative. Penultimate adversarial signal is at most
$3.05\times$ its own permutation noise (below $1\times$ in 16/24 cells); the KM signal sits
$7\times10^{2}$--$1.7\times10^{7}$ above its noise floor. Worst-case tail on ResNet-152
(min signal / max noise): 1.14$\times$ (DeepFool); the ResNet-152 floor is the
\emph{wide-face} permutation subgroup and its residual is float32 accumulation, not exact invariance.
DenseNet-121/GoogLeNet have no Step-A run by design (covered by teleportation); penultimate
per-pair quantiles are unrecoverable from stored summaries.
\clarify{In panel~(b) the bar in ``$x \mid y$'' is a separator, not a division: $x$ is the
knowledge matrix's ratio and $y$ the penultimate features' own, each formed in its own units.
Both cells are \emph{mean adversarial signal} $\div$ \emph{mean permutation noise}, the two
scales defined in Section~\ref{sec:setup}, so a value above $1$ says the representation
registers the attack more strongly than it registers a function-preserving relabelling of
neurons. Worked example (ResNet-152/FGSM): $83.98/0.1157=726$ and $7.600/15.05=0.51$.}}
\label{tab:signal-relative}
\end{table}


\subsection{Reproducing this}
\label{app:teleport-repro}

\texttt{scripts/verify\_teleportation\_exactness.py} runs the precision-scaling
comparison of Table~\ref{tab:teleport-exactness} end to end on CPU for all three
architectures; it applies the three de-downcasting patches described above,
teleports with a fixed seed at each precision, and reports the ratio against the
$\varepsilon_{32}/\varepsilon_{64}$ prediction. It needs no GPU and no cluster
allocation, and runs in a few minutes per architecture.
